\documentclass[11pt]{article}

% Anonymous review version
\usepackage[review]{acl}

% Standard package includes
\usepackage{times}
\usepackage{latexsym}
\usepackage[T1]{fontenc}
\usepackage[utf8]{inputenc}
\usepackage{microtype}
\usepackage{inconsolata}
\usepackage{graphicx}
\usepackage{booktabs}
\usepackage{multirow}
\usepackage{amsmath}
\usepackage{xcolor}
\usepackage{subcaption}
\usepackage{enumitem}

\title{SciFig-Eval: A Multi-Dimensional Benchmark for Evaluating Scientific Figure Understanding in Vision-Language Models}

\author{Anonymous ACL submission}

\begin{document}
\maketitle

\begin{abstract}
We present \textsc{SciFig-Eval}, a benchmark for evaluating vision-language models (VLMs) on scientific figure understanding across two complementary axes, open-ended description and targeted reasoning, each under standard and adversarial conditions. We evaluate 8 VLMs on 250 English scientific figures using checklist-based MQM scoring validated against human judgment (Spearman $\rho = 1.0$, $p < 0.001$), 6 image transforms, psychology-informed hallucination probes grounded in presupposition embedding and anchoring bias, and caption bias resistance testing. We introduce the Admittance-Resistance-Inductance (A-R-I) framework, which decomposes behavioral reliability into independently measurable dimensions through selective blur experiments. Our results reveal that quality and reliability rankings diverge at consequential points. GPT-5.2 ranks first on description quality (MQM 91.6) but admits visual limitations only 6\% of the time, while Gemini 3.1 Pro admits 90\% and leads on resistance (0.91). Phi-4 follows poisoned captions 95\% of the time despite producing coherent descriptions. These divergences, invisible to quality-only benchmarks, demonstrate that current evaluation practices provide an incomplete picture of model readiness for scientific deployment.
\end{abstract}

%!TEX root = ../main.tex
\section{Introduction}
\label{sec:introduction}

The model that writes the best scientific figure descriptions is also the one that most confidently fabricates information it cannot see.
Vision-language models are rapidly entering scientific workflows, from literature review assistants that summarize chart findings to accessibility tools that narrate figures for visually impaired researchers.
When a model misreads a trend line or invents a data point, the downstream consequences propagate silently: a fabricated value becomes a cited fact, a hallucinated comparison shapes a research conclusion.
The stakes demand not only that models describe figures accurately, but that they behave reliably when faced with ambiguity, missing information, or misleading context.
Yet current evaluation practices assess only the first of these requirements, leaving the second entirely unmeasured.

Existing benchmarks for scientific figure understanding evaluate models exclusively through question answering on clean images.
FigureQA~\citep{kahou2018figureqa} and PlotQA~\citep{methani2020plotqa} test binary and numerical reasoning over synthetic or semi-synthetic charts.
ChartQA~\citep{masry2022chartqa} advances to human-authored questions on real-world figures, and ChartBench~\citep{xu2024chartbench} broadens chart type coverage.
These benchmarks measure whether a model produces a correct answer, but none ask whether a model admits when it cannot see, resists false premises embedded in questions, or distinguishes its own visual evidence from authoritative-sounding but incorrect textual context.
A model can score highly on ChartQA while confidently fabricating information about chart elements that do not exist, because no existing benchmark tests for this failure mode.

We present \textsc{SciFig-Eval}, a diagnostic evaluation framework that measures both description quality and behavioral reliability on scientific figures.
The framework is organized around two evaluation axes (open-ended description and targeted reasoning) crossed with two conditions (standard and adversarial), yielding a $2\!\times\!2$ evaluation matrix that structures every experiment.
Adversarial probes are grounded in cognitive science: \emph{presupposition embedding}~\citep{loftus1975leading} tests whether definite articles trick models into accepting non-existent chart elements, \emph{anchoring bias}~\citep{tversky1974judgment} plants plausible but incorrect values, and \emph{sycophantic agreement}~\citep{sharma2024sycophancy} measures whether models echo false captions over their own visual perception.
These probes feed into the \emph{A-R-I behavioral framework}, which decomposes model behavior into three independently measurable dimensions: Admittance (does the model acknowledge visual limitations?), Resistance (does it push back on false premises?), and Inductance (can it infer missing information from context?).
We evaluate 8 models across 250 figures and over 11,000 evaluation instances (Figure~\ref{fig:overview}).

Quality rankings diverge from reliability rankings at precisely the points that matter for deployment.
GPT-5.2 leads on description quality with an MQM score of 91.6, but admits visual limitations only 6\% of the time when asked about elements it cannot see, fabricating answers for the remaining 94\%.
Gemini 3.1 Pro, ranking second on quality at 90.2, admits uncertainty 90\% of the time and achieves the highest resistance to hallucination probes (0.91).
The divergence extends to caption dependency: Phi-4 follows poisoned captions over its own visual evidence 95\% of the time ($R = 0.05$), while presupposition embedding, the technique borrowed from eyewitness testimony research, catches models that successfully resist explicit false values but fail against implicit assumptions about non-existent elements.
A benchmark reporting only quality scores would rank GPT-5.2 and Gemini as near-equivalent, missing the fact that one fabricates 98\% of the time while the other acknowledges uncertainty.

Our contributions are as follows:
\begin{enumerate}
    \item A multi-dimensional evaluation framework for scientific figure understanding that pairs standard quality assessment with psychology-informed adversarial probes across two evaluation axes and two conditions (\S\ref{sec:framework}).
    \item The A-R-I behavioral framework, which decomposes model behavior into Admittance, Resistance, and Inductance, providing independently measurable dimensions of reliability under uncertainty (\S\ref{sec:ari}).
    \item An empirical demonstration that description quality and behavioral reliability are distinct dimensions of VLM competence, with model rankings diverging substantially across these dimensions (\S\ref{sec:results}).
    \item A full benchmark release including the annotated dataset, adversarial probes, evaluation code, and an interactive review dashboard.
\end{enumerate}


%!TEX root = ../main.tex
\section{Related Work}
\label{sec:related_work}

\subsection{Scientific Figure Benchmarks}

Chart understanding benchmarks have progressed from synthetic binary questions \citep{kahou2018figureqa} through numerical reasoning over real-world plots \citep{methani2020plotqa,masry2022chartqa} to complex visual reasoning across diverse chart types \citep{xu2024chartbench}. \citet{li2023scigraphqa} targeted scientific graphs specifically, generating multi-turn QA from arXiv papers. A comprehensive comparison of these and additional benchmarks appears in Appendix~\ref{sec:benchmark-comparison}.

These benchmarks evaluate perception and reasoning. They ask whether a model can extract a value, identify a trend, or answer a factual question about a clean image. None evaluate \emph{behaviour}, whether a model fabricates information it cannot see, echoes false premises, or acknowledges uncertainty when visual evidence is degraded. \textsc{SciFig-Eval} fills this gap by adding a behavioural evaluation axis, measuring not only what a model says about a figure but how it acts when challenged with partial visibility and misleading context. The perception-reasoning-behaviour distinction structures every experiment in this paper and reveals that models ranking identically on the first two dimensions diverge dramatically on the third (\S\ref{sec:results}).

\subsection{VLM Hallucination and Reliability}

Hallucination research in vision-language models has focused on natural images. \citet{rohrbach2018object} introduced the CHAIR metric for object hallucination in image captions. \citet{li2023pope} proposed POPE, converting hallucination detection into a polling-based binary task that probes co-occurrence biases. \citet{guan2024hallusionbench} designed HallusionBench with expert-crafted visual illusions, and \citet{sun2024aligning} introduced MMHal-Bench to penalise hallucination across eight question categories. \citet{bai2024hallucination} provide a comprehensive survey categorising hallucination sources.

Scientific figures differ from natural images in ways that make existing hallucination frameworks insufficient. Structured layouts, quantitative data encoded in visual channels, and domain conventions demand precise extraction rather than approximate recognition. A model hallucinating a non-existent bar in a chart carries different consequences from a model inventing a background object in a photograph. More fundamentally, no existing benchmark tests whether a model \emph{admits} when it cannot see.

Our behavioural probes draw on cognitive science rather than treating hallucination as a monolithic category. \emph{Presupposition embedding}~\citep{loftus1975leading} uses definite articles to make models accept non-existent chart elements. \emph{Anchoring bias}~\citep{tversky1974judgment} plants plausible but incorrect values that models fail to challenge. The \emph{cooperative principle}~\citep{grice1975logic} exploits models' assumption that questions are well-formed, producing answers rather than refusals. \emph{Sycophantic agreement}~\citep{sharma2024sycophancy} captures the tendency to echo provided captions over direct visual evidence. Together, these probes operationalise the behavioural dimension that existing benchmarks leave unmeasured.

\subsection{Evaluation Methodology}

\textsc{SciFig-Eval} adapts the Multidimensional Quality Metrics (MQM) framework~\citep{lommel2014multidimensional}, originally developed for translation quality assessment and adopted as the standard for WMT human evaluation~\citep{freitag2021experts}. MQM decomposes quality into fine-grained error categories with severity weightings, enabling diagnostic evaluation rather than holistic scoring. We extend MQM to figure descriptions by introducing \emph{binding verification}, which checks not only whether described elements exist in the figure but whether stated relationships between elements (e.g., ``the blue bar exceeds the red bar'') are grounded in the visual evidence.

For scalable evaluation, we follow the LLM-as-judge paradigm validated by \citet{zheng2023judging}, who demonstrated strong agreement between LLM judges and human preferences on MT-Bench. Our judges apply checklist-based scoring in which each figure is annotated with verifiable claims, and the judge evaluates whether the description correctly addresses each claim. This structured approach mitigates the subjectivity of holistic ratings while enabling evaluation of eight models across multiple conditions.

Existing benchmarks measure perception and reasoning \emph{or} detect hallucinations. \textsc{SciFig-Eval} measures all three dimensions and finds that model rankings diverge substantially across them (\S\ref{sec:results}).


%!TEX root = ../main.tex
\section{Evaluation Framework}
\label{sec:framework}

\textsc{SciFig-Eval} evaluates vision-language models along three dimensions of increasing demand: \emph{perception} (can the model accurately describe a scientific figure?), \emph{reasoning} (can it answer analytical questions about what it sees?), and \emph{behaviour} (how does it act under uncertainty and deception?). A model that excels at the first two may still fail catastrophically at the third. The framework is organized as a $2\!\times\!2$ evaluation matrix (Table~\ref{tab:eval-matrix}) crossing two evaluation modes (open-ended description and targeted reasoning) with two conditions (standard and adversarial).

\begin{table}[t]
\centering
\small
\begin{tabular}{@{}lp{3.1cm}p{3.1cm}@{}}
\toprule
 & \textbf{Standard} & \textbf{Adversarial} \\
\midrule
\textbf{Description} & Baseline MQM & Caption bias, passive admittance / inductance \\
\textbf{Reasoning}    & Capability Qs & Resistance probes, active admittance / inductance \\
\bottomrule
\end{tabular}
\caption{The $2\!\times\!2$ evaluation matrix. Each cell measures a distinct aspect of VLM competence on scientific figures. Standard conditions assess perception and reasoning; adversarial conditions expose behaviour under uncertainty and deception.}
\label{tab:eval-matrix}
\end{table}

Standard conditions measure perception through checklist-based MQM scoring and reasoning through targeted capability questions (counting, computation, comparison, pattern analysis). Adversarial conditions probe behaviour by presenting models with two types of challenge. \emph{Uncertainty} tests whether a model acknowledges what it cannot see, measured through selective blur of chart elements. \emph{Deception} tests whether a model resists false information, measured through poisoned captions and presupposition-laden questions. These behavioural probes feed into the A-R-I framework (\S\ref{sec:ari}), which decomposes behaviour into three independently measurable dimensions.

\subsection{Checklist-Based MQM Evaluation}
\label{sec:mqm}

Perception is measured through a chart-type-specific variant of the Multidimensional Quality Metrics (MQM) framework. Each chart type has a dedicated checklist covering structural identification (chart type, orientation, grouping), axis and legend accuracy, numerical correctness, and visual element description. Bar charts receive 14~items, line plots 15, and pie charts 11. Each item is classified as \emph{Major} or \emph{Minor} based on its impact on figure comprehension.

An LLM judge (GPT-4o) assesses each checklist item on two axes: \emph{coverage} (complete, partial, or missing) and \emph{correctness} (correct, partial, wrong, or n/a), then scans for violations of seven global constraints. A rule-based engine maps judge output to penalties across three dimensions (Accuracy, Completeness, and Clarity) with severity-dependent weights (full weight table in Appendix~\ref{sec:mqm-details}). The final MQM score normalizes total penalties against a severity-adjusted maximum:
\begin{equation}
\text{MQM} = \max\!\Bigl(0,\; 100 - \frac{\sum_{i} w_i}{D} \times 100\Bigr)
\label{eq:mqm}
\end{equation}

\noindent where $w_i$ is the weight of penalty~$i$ and $D$ is the maximum possible penalty given the checklist composition.

Two mechanisms guard against scoring artifacts. \emph{Binding verification} cross-checks element-level value-label-colour correspondences, escalating any mismatch to a Major error. A \emph{deduplication engine} collapses overlapping penalties from global constraints and checklist items. To validate automated scoring, two expert annotators independently ranked all eight models on a 30-figure subset. The automated MQM rankings achieve a Spearman correlation of $\rho = 1.0$ ($p < 0.001$, $n = 8$ models) with the human consensus ranking, confirming that the checklist-based approach preserves human judgment at scale.

\subsection{Caption Bias Probes}
\label{sec:caption-bias}

Caption bias probes test whether a model trusts its own visual perception over misleading textual context, targeting the deception axis of behaviour. For each figure, we generate a modified caption containing exactly 2--3 embedded false claims, designed according to psychology-informed principles. The \emph{70/30 rule} ensures that roughly 70\% of the modified caption is verifiably correct, establishing trust before introducing poisoned claims. False numerical values occupy an \emph{anchoring sweet spot} of 20--40\% deviation from actual values, close enough to be plausible yet clearly falsifiable by visual inspection \citep{tversky1974judgment}. Claims target secondary details rather than dominant features (\emph{peripheral misinformation}), and framing employs \emph{loaded verbs} that subtly mischaracterize trends (e.g., ``declined'' vs.\ ``stabilized'').

Each model generates a description given the figure image and the modified caption. An LLM judge evaluates each false claim independently through a randomized A/B design, receiving the claim and the ground-truth reality in randomized order to eliminate position bias. The \emph{caption bias resistance score} is:
\begin{equation}
R_{\text{caption}} = \frac{n_{\text{image}}}{n_{\text{image}} + n_{\text{caption}}}
\label{eq:caption-bias}
\end{equation}

\noindent where $n_{\text{image}}$ is the number of claims where the model followed visual evidence and $n_{\text{caption}}$ is the number where it echoed the false caption.

\subsection{Resistance Probes}
\label{sec:resistance}

Resistance probes test whether models push back on false premises rather than fabricating compliant answers. We design three probe types grounded in cognitive science, each exploiting a different mechanism of deception.

\paragraph{Inexist probes} use \emph{presupposition embedding} \citep{loftus1975leading}. Definite articles (``the error bars,'' ``the regression line'') presuppose the existence of chart elements that are absent. The false element is embedded in a subordinate clause, directing attention to an analytical question \emph{about} the element rather than \emph{whether} it exists. This exploits co-occurrence priors for the chart type.

\paragraph{Contra probes} exploit \emph{anchoring bias} \citep{tversky1974judgment}. A specific but incorrect numerical value (20--30\% off) is embedded as a premise, and the model is asked to build upon it. The model must first detect the false anchor, then correct it, a task substantially harder than reading a value directly.

\paragraph{Unanswerable probes} leverage \emph{conversational implicature} \citep{grice1975logic}. Domain-appropriate questions request $p$-values, sample sizes, or projections beyond the plotted range. These sound like legitimate analytical follow-ups but cannot be answered from the chart alone. Refusal is costly because the question sounds helpful and reasonable.

Each probe response is scored on a three-point rubric: 1.0~(correctly resists), 0.5~(hedges but partially complies), and 0.0~(fully accepts the false premise or fabricates an answer). The aggregate \emph{resistance score} averages across all probes per model.

\subsection{The A-R-I Behavioural Framework}
\label{sec:ari}

The evaluation components above assess perception (MQM), reasoning (capability questions), and specific behavioural probes (caption bias, resistance). We unify the behavioural dimension through the \emph{A-R-I framework}, which decomposes model behaviour into three independently measurable dimensions. Together, they answer a question that quality scores alone cannot: when the input is degraded or deceptive, does the model behave responsibly?

\paragraph{Admittance} measures behaviour under uncertainty. When chart elements are selectively blurred, does the model acknowledge what it cannot see, or does it fabricate confidently? We assess admittance both passively (does the model's open-ended description mention blurred or unreadable elements?) and actively (when asked a targeted question about a blurred element, does it admit uncertainty?). Active admittance evaluation uses a judge that independently scores two binary dimensions: whether the model \emph{admits} uncertainty and whether it \emph{fabricates} a specific answer. A model can do both simultaneously (e.g., ``the label is obscured but appears to be X''), and this decomposition captures important behavioural nuance.

\paragraph{Resistance} measures behaviour under deception. This dimension aggregates the resistance probes (\S\ref{sec:resistance}) and caption bias probes (\S\ref{sec:caption-bias}), capturing whether the model prioritizes its own visual evidence over authoritative-sounding but false textual context. A model with high resistance pushes back on false premises; one with low resistance complies with whatever it is told.

\paragraph{Inductance} measures reasoning under uncertainty. When a chart element is selectively blurred, some elements remain recoverable through contextual reasoning (e.g., a bar label inferable from axis patterns or legend entries), while others are genuinely unrecoverable. We assess inductance through the same selective blur methodology used for admittance, but focus on the \emph{correctness} of fabricated answers for inferable elements. When a model fabricates an answer for an element that could be inferred from context, correctness validates whether genuine reasoning occurred or the output was a lucky guess. This distinction is empirically validated by a large gap in fabrication correctness between inferable and non-inferable elements (\S\ref{sec:results}), confirming that inductance captures real reasoning rather than noise.

\paragraph{Independence of dimensions.} The three dimensions vary independently across models. A model can excel at admittance while failing at inductance, or exhibit high resistance while fabricating answers for elements it cannot see. As we demonstrate in \S\ref{sec:results}, models that rank highest on perception do not necessarily rank highest on behaviour, and the gap between the two can be dramatic.


\begin{table*}[t]
\centering
\small
\begin{tabular}{@{}lrrrrrrrrrr@{}}
\toprule
\textbf{Model} & \textbf{Base} & \textbf{Orig.} & \textbf{Noise} & \textbf{Rot.} & \textbf{LowC} & \textbf{InPap} & \textbf{IPBlr} & \textbf{CapB} & \textbf{AdmB} & \textbf{IndB} \\
\midrule
GPT-5.2 & 91.6$_{\scriptstyle 90.4}^{\scriptstyle 92.8}$ & 89.6$_{\scriptstyle 87.1}^{\scriptstyle 91.8}$ & 91.3$_{\scriptstyle 89.1}^{\scriptstyle 93.3}$ & 70.2$_{\scriptstyle 66.5}^{\scriptstyle 73.9}$ & 87.0$_{\scriptstyle 84.1}^{\scriptstyle 89.6}$ & 77.8$_{\scriptstyle 74.2}^{\scriptstyle 81.3}$ & 13.3$_{\scriptstyle 9.2}^{\scriptstyle 17.8}$ & 91.3$_{\scriptstyle 89.2}^{\scriptstyle 93.3}$ & 81.7$_{\scriptstyle 77.4}^{\scriptstyle 85.6}$ & 88.9$_{\scriptstyle 84.0}^{\scriptstyle 93.2}$ \\
Gemini & 90.2$_{\scriptstyle 88.9}^{\scriptstyle 91.4}$ & 89.3$_{\scriptstyle 87.2}^{\scriptstyle 91.3}$ & 90.2$_{\scriptstyle 88.2}^{\scriptstyle 92.1}$ & 72.0$_{\scriptstyle 68.6}^{\scriptstyle 75.3}$ & 85.1$_{\scriptstyle 82.7}^{\scriptstyle 87.4}$ & 83.8$_{\scriptstyle 81.1}^{\scriptstyle 86.3}$ & 8.4$_{\scriptstyle 4.9}^{\scriptstyle 12.1}$ & 91.3$_{\scriptstyle 89.4}^{\scriptstyle 93.2}$ & 84.6$_{\scriptstyle 81.5}^{\scriptstyle 87.5}$ & 86.8$_{\scriptstyle 82.7}^{\scriptstyle 90.4}$ \\
Llama 4 & 81.4$_{\scriptstyle 79.5}^{\scriptstyle 83.3}$ & 82.0$_{\scriptstyle 78.9}^{\scriptstyle 84.9}$ & 81.0$_{\scriptstyle 77.9}^{\scriptstyle 84.0}$ & 60.6$_{\scriptstyle 56.3}^{\scriptstyle 64.7}$ & 75.9$_{\scriptstyle 72.7}^{\scriptstyle 78.8}$ & 63.0$_{\scriptstyle 58.9}^{\scriptstyle 66.9}$ & 30.9$_{\scriptstyle 26.7}^{\scriptstyle 35.2}$ & 84.8$_{\scriptstyle 81.7}^{\scriptstyle 87.5}$ & 70.5$_{\scriptstyle 65.7}^{\scriptstyle 75.1}$ & 78.9$_{\scriptstyle 74.2}^{\scriptstyle 83.5}$ \\
Qwen-235B & 80.8$_{\scriptstyle 78.9}^{\scriptstyle 82.7}$ & 82.6$_{\scriptstyle 79.8}^{\scriptstyle 85.3}$ & 82.1$_{\scriptstyle 79.6}^{\scriptstyle 84.7}$ & 65.3$_{\scriptstyle 62.0}^{\scriptstyle 68.6}$ & 77.8$_{\scriptstyle 74.7}^{\scriptstyle 80.6}$ & 77.2$_{\scriptstyle 73.6}^{\scriptstyle 80.7}$ & 22.9$_{\scriptstyle 18.3}^{\scriptstyle 27.7}$ & 83.4$_{\scriptstyle 80.9}^{\scriptstyle 85.7}$ & 71.0$_{\scriptstyle 65.2}^{\scriptstyle 76.7}$ & 79.0$_{\scriptstyle 74.5}^{\scriptstyle 83.1}$ \\
Qwen-8B & 78.9$_{\scriptstyle 76.8}^{\scriptstyle 81.1}$ & 78.5$_{\scriptstyle 75.0}^{\scriptstyle 81.8}$ & 80.1$_{\scriptstyle 76.8}^{\scriptstyle 83.2}$ & 61.9$_{\scriptstyle 58.2}^{\scriptstyle 65.4}$ & 76.0$_{\scriptstyle 72.4}^{\scriptstyle 79.5}$ & 72.9$_{\scriptstyle 68.9}^{\scriptstyle 76.9}$ & 25.8$_{\scriptstyle 21.3}^{\scriptstyle 30.4}$ & 83.0$_{\scriptstyle 80.1}^{\scriptstyle 85.8}$ & 71.5$_{\scriptstyle 66.1}^{\scriptstyle 76.8}$ & 74.8$_{\scriptstyle 68.9}^{\scriptstyle 80.2}$ \\
Qwen-30B & 74.4$_{\scriptstyle 72.1}^{\scriptstyle 76.7}$ & 77.3$_{\scriptstyle 73.6}^{\scriptstyle 80.7}$ & 74.3$_{\scriptstyle 70.6}^{\scriptstyle 78.0}$ & 58.2$_{\scriptstyle 54.1}^{\scriptstyle 62.1}$ & 75.7$_{\scriptstyle 72.5}^{\scriptstyle 78.9}$ & 68.3$_{\scriptstyle 63.7}^{\scriptstyle 72.5}$ & 24.3$_{\scriptstyle 19.9}^{\scriptstyle 28.7}$ & 78.3$_{\scriptstyle 74.7}^{\scriptstyle 81.6}$ & 67.5$_{\scriptstyle 61.5}^{\scriptstyle 73.1}$ & 74.7$_{\scriptstyle 69.1}^{\scriptstyle 80.0}$ \\
Gemma & 69.1$_{\scriptstyle 66.6}^{\scriptstyle 71.6}$ & 60.9$_{\scriptstyle 56.7}^{\scriptstyle 65.1}$ & 62.2$_{\scriptstyle 58.0}^{\scriptstyle 66.5}$ & 49.8$_{\scriptstyle 45.7}^{\scriptstyle 54.1}$ & 61.5$_{\scriptstyle 57.8}^{\scriptstyle 65.3}$ & 35.4$_{\scriptstyle 31.5}^{\scriptstyle 39.6}$ & 25.3$_{\scriptstyle 21.6}^{\scriptstyle 29.1}$ & 70.9$_{\scriptstyle 65.6}^{\scriptstyle 75.8}$ & 58.3$_{\scriptstyle 53.7}^{\scriptstyle 63.0}$ & 59.6$_{\scriptstyle 52.4}^{\scriptstyle 66.6}$ \\
Phi-4 & 62.2$_{\scriptstyle 59.3}^{\scriptstyle 64.9}$ & 59.5$_{\scriptstyle 54.2}^{\scriptstyle 64.6}$ & 61.3$_{\scriptstyle 57.1}^{\scriptstyle 65.3}$ & 43.7$_{\scriptstyle 39.4}^{\scriptstyle 48.0}$ & 64.7$_{\scriptstyle 60.5}^{\scriptstyle 68.9}$ & 31.3$_{\scriptstyle 26.4}^{\scriptstyle 36.3}$ & 17.7$_{\scriptstyle 14.4}^{\scriptstyle 21.1}$ & 61.7$_{\scriptstyle 56.1}^{\scriptstyle 67.1}$ & 56.6$_{\scriptstyle 50.3}^{\scriptstyle 62.3}$ & 59.4$_{\scriptstyle 52.6}^{\scriptstyle 66.1}$ \\
\bottomrule
\end{tabular}
\caption{Description quality (MQM scores, 0--100) across conditions. Base = baseline with caption (250 figures). Transforms evaluated on 100-figure subset. Values show mean with 95\% bootstrap CI (subscript/superscript). Higher is better.}
\label{tab:mqm-results}
\end{table*}
\begin{table*}[t]
\centering
\small
\setlength{\tabcolsep}{4pt}
\begin{tabular}{@{}l rrrr rrrr rr rr@{}}
\toprule
& \multicolumn{4}{c}{\textbf{Capability (\%)}} & \multicolumn{4}{c}{\textbf{Resistance}} & \multicolumn{2}{c}{\textbf{Admittance (\%)}} & \multicolumn{2}{c}{\textbf{Inductance (\%)}} \\
\cmidrule(lr){2-5} \cmidrule(lr){6-9} \cmidrule(lr){10-11} \cmidrule(lr){12-13}
\textbf{Model} & \textbf{Cnt} & \textbf{Cmp} & \textbf{Cmpr} & \textbf{Pat} & \textbf{Inex} & \textbf{Cont} & \textbf{Unan} & \textbf{CapB} & \textbf{Act} & \textbf{Pas} & \textbf{Act} & \textbf{Pas} \\
\midrule
\mdot{clrGemini}\textbf{Gemini} & \textbf{89.2} & 79.4 & \textbf{89.6} & 70.0 & \textbf{.88} & \textbf{.91} & \textbf{.95} & .89 & \textbf{71} & \textbf{59} & \textbf{66} & \underline{73} \\
\mdot{clrGPT}\underline{GPT-5.2} & \underline{76.1} & \textbf{82.8} & \underline{77.9} & \textbf{72.0} & \underline{.77} & \underline{.75} & .92 & \textbf{.89} & 8 & \underline{23} & \underline{59} & \textbf{77} \\
\mdot{clrLlama}Llama 4 & 45.6 & \underline{53.4} & 37.2 & 48.0 & .63 & .76 & \underline{.94} & .74 & \underline{19} & 5 & 34 & 58 \\
\mdot{clrQwen235}Qwen-235B & 65.2 & 63.7 & 50.0 & 52.0 & .67 & .64 & \underline{.94} & .54 & 15 & 13 & 29 & 58 \\
\mdot{clrQwen8}Qwen-8B & 47.8 & 52.8 & 45.4 & 50.0 & .40 & .44 & .88 & .43 & 7 & 3 & 22 & 52 \\
\mdot{clrQwen30}Qwen-30B & 43.5 & 45.9 & 31.4 & 30.0 & .23 & .37 & .73 & .30 & 7 & 0 & 24 & 58 \\
\mdot{clrGemma}Gemma & 15.2 & 29.4 & 18.6 & \underline{40.0} & .17 & .24 & .93 & .38 & 8 & 2 & 14 & 41 \\
\mdot{clrPhi}Phi-4 & 13.0 & 6.2 & 3.5 & 16.0 & .04 & .04 & .56 & .05 & 5 & 2 & 15 & 35 \\
\bottomrule
\end{tabular}
\caption{Unified behavioural evaluation. \textbf{Capability} shows accuracy (\%) on four question types (Cnt = counting, Cmp = computation, Cmpr = comparison, Pat = pattern analysis; 250 figures, averaged across two judges). \textbf{Resistance} shows scores (0--1) for three resistance probe types (Inex = inexist, Cont = contra, Unan = unanswerable; 250 figures) and caption bias resistance (CapB; 100 figures). \textbf{Admittance} shows the percentage of probes where the model acknowledged visual uncertainty (Act = active targeted question, Pas = passive open-ended description; 228 figures). \textbf{Inductance} shows the percentage of fabricated answers that were correct for context-inferable elements (Act = active, Pas = passive; 215 figures). Capability, Admittance, and Inductance columns report percentages; Resistance columns report scores on a 0--1 scale. \textbf{Bold} = best per column, \underline{underline} = second best.}
\label{tab:behavioral}
\end{table*}


%!TEX root = ../main.tex
\section{Experiments and Results}
\label{sec:results}

We evaluate eight vision-language models across 250 scientific figures and over 14 experimental conditions, yielding more than 11,000 individual evaluation instances. The central finding is that model rankings under perception, reasoning, and behavioural evaluation are not the same ranking. GPT-5.2 produces the highest-quality descriptions in our benchmark (MQM 91.6) yet fabricates answers for elements it cannot see 94\% of the time, with nothing in its output to distinguish the fabrication from genuine analysis.

\subsection{Models}
\label{sec:models}

We select eight models spanning commercial and open-weight families, dense and mixture-of-experts (MoE) architectures, and a range of effective parameter counts.
\textbf{Commercial models}: GPT-5.2 (Azure, closed-weight), Gemini 3.1 Pro (closed-weight), and Phi-4-Multimodal (small, closed-weight).
\textbf{Open-weight MoE models}: Llama 4 Maverick (17B params, 128 experts), Qwen3-VL-235B-A22B (235B total, 22B active), and Qwen3-VL-30B-A3B (30B total, 3B active).
\textbf{Open-weight dense models}: Qwen3-VL-8B and Gemma3-27B-IT.
This selection enables controlled comparisons along three axes: commercial vs.\ open-weight deployment, MoE vs.\ dense architecture at similar active parameter counts (Qwen-235B at 22B active vs.\ Gemma at 27B dense), and scaling within a single family (the three Qwen variants from 3B to 22B active parameters). All models received identical prompts and were evaluated at temperature~0 with GPT-4o as the automated judge.

\subsection{Perception}
\label{sec:perception-results}

\paragraph{Baseline description quality.}
GPT-5.2 leads with a baseline MQM of 91.6 [90.4, 92.8], followed by Gemini at 90.2 [88.9, 91.4] (Table~\ref{tab:mqm-results}). The 1.4-point gap is statistically significant ($p < 0.01$) but practically small (Cliff's $\delta = 0.09$, negligible). Below these two leaders, a clear tier structure emerges. Llama~4, Qwen-235B, and Qwen-8B cluster between 78.9 and 81.4 with overlapping confidence intervals and negligible pairwise effect sizes. Qwen-30B (74.4) and Gemma (69.1) form a third tier. Phi-4 trails at 62.2 [59.3, 64.9], nearly 30 points behind the leaders, with completeness penalties (8.8 points) driving the gap far more than accuracy penalties do for stronger models (GPT-5.2 loses only 0.7 points to completeness).

\paragraph{Transform robustness.}
Rotation is the most damaging transform, causing an average MQM drop of 19.4 points across models, while noise has negligible impact (Table~\ref{tab:mqm-results}). Low contrast reduces scores by a moderate 4--7 points for top-tier models. The most revealing condition is in-paper-blur, where the figure is embedded in a PDF page and the surrounding content is blurred. This condition collapses all models below 31 MQM points. Gemini drops from 90.2 to 8.4. GPT-5.2 drops from 91.6 to 13.3. Both attempt to describe the blurred page context rather than focusing on the embedded figure, a failure with practical implications for document-level figure understanding pipelines.

\paragraph{Selective blur.}
Selectively blurring unrecoverable elements (admittance blur) reduces MQM by roughly 8--10 points for top models, while blurring inferable elements (inductance blur) produces smaller drops of 1--3 points (Table~\ref{tab:mqm-results}, AdmB and IndB columns). This difference validates the admittance-inductance distinction at the perception level. GPT-5.2 scores 88.9 under inductance blur but 81.7 under admittance blur, a 7.2-point gap that reflects the difference between elements a model can reason about and elements that are simply gone.

\paragraph{Caption bias on description quality.}
Poisoned captions have no measurable effect on MQM for resistant models. GPT-5.2 and Gemini both score 91.3 under modified captions, matching their baselines, because they ignore the false claims entirely. Phi-4, already caption-dependent, scores 61.7 under the same condition, consistent with its near-zero resistance ($R = 0.05$).

\subsection{Reasoning}
\label{sec:reasoning-results}

Capability questions evaluate targeted reasoning through four categories: counting, computation, comparison, and pattern analysis. Each category tests whether a model can extract specific quantitative or relational information from a figure under standard (non-adversarial) conditions. Full results appear in Table~\ref{tab:capability}.\footnote{Capability evaluation is currently being finalized across all eight models; we report preliminary results in the appendix and will include the complete table in the final version.}

\subsection{Behaviour}
\label{sec:behaviour-results}

Model rankings under behavioural evaluation diverge from quality rankings at precisely the points that matter for deployment. GPT-5.2 ranks first on description quality but falls to fourth on active admittance, while Gemini leads on nearly every behavioural dimension despite placing second on quality (Table~\ref{tab:behavioral}).

\paragraph{Resistance to hallucination probes.}
Gemini achieves the highest overall resistance at 0.91 [0.89, 0.93], followed by GPT-5.2 at 0.81 [0.79, 0.84]. Phi-4 anchors the bottom at 0.21 [0.18, 0.24]. The gap between Gemini and GPT-5.2 is statistically significant ($p < 0.001$, $n = 750$ probe instances), confirming that resistance is not simply a correlate of overall model capability.

Probe type difficulty follows a clear gradient. Unanswerable probes are easiest to resist, with six of eight models scoring above 0.88, because the question format permits straightforward refusal. Inexist probes, grounded in presupposition embedding \citep{loftus1975leading}, are hardest, with scores ranging from 0.04 (Phi-4) to 0.88 (Gemini). This 22$\times$ gap reflects the power of definite articles to bypass model vigilance. A model that can refuse an explicitly unanswerable question still accepts the existence of chart elements that were never there. Contra probes fall between these extremes, with anchoring bias \citep{tversky1974judgment} proving moderately effective at suppressing model pushback.

\paragraph{Caption bias resistance.}
Models fall along a caption dependency spectrum from near-total visual independence to near-total textual dependency (Table~\ref{tab:behavioral}, Cap.\ Bias~R). Gemini and GPT-5.2 both achieve 0.89 [0.85, 0.93], with overlapping CIs and a non-significant difference ($p = 0.44$). Llama~4 provides moderate resistance at 0.74. Phi-4 exhibits near-zero resistance at 0.05 [0.02, 0.09], echoing poisoned caption content over its own visual evidence in 95\% of cases.

Caption dependency tracks active parameter count within the Qwen family: 0.43 (8B) $\rightarrow$ 0.30 (30B, 3B active) $\rightarrow$ 0.54 (235B, 22B active). The non-monotonic pattern suggests that active parameter count, rather than total parameter count, predicts caption independence for MoE architectures.

\paragraph{Active admittance and inductance.}
Gemini is the only model that consistently acknowledges visual limitations, admitting uncertainty in 90\% of cases [82\%, 98\%] when asked about selectively blurred elements (Table~\ref{tab:behavioral}). No other model exceeds 22\%. GPT-5.2, the top-ranked model on description quality, admits visual limitations only 6\% of the time while fabricating answers 98\% of the time. This ``confident fabricator'' profile, high descriptive fluency coexisting with near-total unwillingness to acknowledge uncertainty, represents a concrete safety risk for scientific deployment.

The inductance dimension validates the A-R-I framework empirically. When models fabricate answers for \emph{inferable} elements (inductance probes), correctness ranges from 21\% to 80\% across models (Table~\ref{tab:behavioral}). When they fabricate for \emph{unrecoverable} elements (admittance probes), correctness drops to 0--14\%, consistent with lucky guessing. GPT-5.2 achieves 74\% inductance correctness ($n = 47$ fabrications) versus 14\% admittance correctness ($n = 49$), a 60-percentage-point gap that confirms models perform genuine contextual reasoning when context permits rather than generating plausible-sounding noise.

\paragraph{Passive admittance.}
Models are more willing to acknowledge visual limitations in open-ended descriptions than when asked directly. GPT-5.2 mentions blur or illegibility for 26\% of admittance-blurred elements in its descriptions but admits uncertainty for only 6\% when asked a targeted question (Table~\ref{tab:behavioral}). Gemini shows the same pattern at a higher baseline (74\% passive vs.\ 90\% active). This discrepancy suggests that targeted questions create social pressure to provide a definitive answer, consistent with Grice's cooperative principle \citep{grice1975logic}, while open-ended descriptions provide more natural opportunities for hedging.

\subsection{Ablation: Probe Designer Independence}
\label{sec:ablation}

A potential concern with LLM-generated probes is circularity. If the same model family designs the probes and evaluates the responses, results may reflect shared biases rather than genuine model behaviour. We compare probes designed by GPT-4o (our default) against probes designed by Mistral Large, evaluating both on the same 50-figure subset with GPT-5.2 as the target model. Caption bias resistance is nearly identical: 0.89 (GPT-4o probes) vs.\ 0.89 (Mistral probes), with a negligible effect size ($\delta = -0.01$, $p = 0.49$). Hallucination resistance shows a small difference (0.80 vs.\ 0.86, $\delta = -0.16$) whose 95\% CI includes zero. These results indicate that the behavioural patterns we observe are properties of the evaluated models, not artifacts of the probe designer.


%!TEX root = ../main.tex
\section{Analysis}
\label{sec:analysis}

Description quality and behavioral reliability, while positively correlated at the population level, diverge at precisely the points that matter for deployment decisions. GPT-5.2 ranks first on MQM (91.6) but fourth on active admittance (6\%), while Gemini ranks second on MQM (90.2) but first on admittance (90\%), resistance (0.91), and caption bias resistance (0.89). The Spearman correlations in Table~\ref{tab:cross-dim} are high overall ($\rho = 0.83$--$0.95$), yet they mask these consequential reversals because the bottom of the ranking is concordant across dimensions. Split-half reliability of $\rho = 0.979$ over 100 random splits (Table~\ref{tab:stability}) confirms that this divergence is stable, not an artifact of sample variance. The practical implication is clear: a benchmark that reports only MQM would rank GPT-5.2 and Gemini as near-equivalent, missing the fact that one fabricates answers for blurred elements 98\% of the time while the other admits uncertainty 90\% of the time.

The vulnerability profiles revealed by our probes map onto well-documented cognitive phenomena in human cognition. Presupposition embedding (inexist probes) is the most effective deception technique across all eight models. Llama~4, for example, resists explicit false values at 0.76 (contra) but drops to 0.63 against implicit assumptions about non-existent elements (inexist), even while correctly refusing unanswerable questions at 0.94. This asymmetry parallels findings from eyewitness testimony research, where definite articles (``the broken headlight'') reliably induce false memories of objects that were never present \citep{loftus1975leading}. Caption bias resistance reveals a separate and unexpected pattern. Phi-4 follows poisoned captions almost entirely ($R = 0.05$), yet Gemma resists at $R = 0.38$ despite scoring 7 MQM points lower (69.1 vs.\ 62.2). If caption dependency were simply a capability limitation, weaker models should be uniformly more susceptible. Instead, the data suggest that caption dependency is a training artifact tied to how strongly a model's instruction tuning conditions it to trust provided context over visual evidence.

The most striking behavioral asymmetry emerges in how models handle visual uncertainty across evaluation modes. In passive descriptions, GPT-5.2 acknowledges blurred elements 26\% of the time, but when directly questioned about those same elements, its admittance drops to 6\%. Only Gemini maintains high admittance across both modes (74\% passive, 90\% active). This pattern suggests that direct questions trigger a ``must answer'' bias, consistent with models trained under RLHF to maximize helpfulness \citep{sharma2023towards}. Open-ended descriptions provide more latitude for natural hedging (``the label appears obscured''), while targeted questions compress the response space toward definitive answers. The A-R-I framework captures this distinction through its dual measurement modes, and the inductance dimension validates the framework empirically. When blurred elements are inferable from context, models answer correctly 21--80\% of the time (GPT-5.2 at 74\%, Gemini at 80\%), but when elements are genuinely unrecoverable, correctness drops to 0--14\% ($n = 50$ per model). This gap confirms that the framework measures real reasoning rather than random fabrication.

These findings are robust to methodological choices. The probe designer ablation (Table~\ref{tab:ablation}) shows that switching from GPT-4o to Mistral Large~3 as the probe generator yields negligible differences in caption bias resistance ($\Delta = 0.002$, $p = 0.49$, Cliff's $\delta = -0.007$) and only a small effect on hallucination resistance ($\Delta = -0.057$, Cliff's $\delta = -0.16$), confirming that our results reflect genuine model behavior rather than idiosyncrasies of the probe generator. Scale validation further supports stability: resistance scores computed on 100 figures closely match those from the full 250-figure set, with maximum model-level deviation of 0.02 points (Table~\ref{tab:stability}).


%!TEX root = ../main.tex
\section{Conclusion}
\label{sec:conclusion}

Quality benchmarks for scientific figure understanding tell an incomplete story. By evaluating eight vision-language models across both description quality and behavioral reliability, \textsc{SciFig-Eval} reveals that these two dimensions are largely independent. The A-R-I framework decomposes model behavior into admittance, resistance, and inductance, exposing failure modes that quality scores alone cannot detect. GPT-5.2 writes the best descriptions yet fabricates answers for blurred elements 98\% of the time, while Gemini, scoring only 1.4 MQM points lower, acknowledges visual limitations in 90\% of those same cases.

These findings carry direct consequences for deploying VLMs in scientific workflows. Selecting a model for automated figure analysis based on quality rankings alone risks embedding silent fabrication into research pipelines. The psychology-informed probes in \textsc{SciFig-Eval} offer a principled toolkit for stress-testing model honesty before deployment, grounded in decades of cognitive science on presupposition embedding, anchoring bias, and cooperative communication. The A-R-I framework itself is not specific to scientific figures and can be applied wherever models must reason under uncertainty, refuse misleading premises, or infer from partial evidence.

Several directions extend this work naturally. Broadening coverage to tables, multi-panel figures, and scientific diagrams would test whether the behavioral patterns we observe generalize across visual genres. Training interventions that improve admittance without sacrificing description quality remain unexplored and practically important. Scaling \textsc{SciFig-Eval} to additional languages would reveal whether the passive-active admittance gap persists across linguistic and cultural contexts. Finally, prompt-based strategies for closing that gap, encouraging models to express uncertainty when directly questioned, offer a low-cost path toward safer scientific deployment.

\section*{Limitations}

Our evaluation focuses on English-language bar charts, line plots, and pie charts. This scope was a deliberate design decision that enabled depth of analysis (over 11,000 evaluation instances from 250 figures), but it means our findings may not transfer to other chart types such as scatter plots, heatmaps, or schematic diagrams, nor to figures drawn from domains outside the computer science and machine learning literature prevalent on arXiv.

The 250-figure dataset, while yielding more than 23,000 model-output instances across all conditions, remains modest in absolute count. We validated stability through split-half reliability ($\rho = 0.979$) and scale analysis showing convergence at 100 figures, yet generalization to broader scientific corpora warrants further investigation. The admittance and inductance conditions were evaluated on 50 and 48 figures respectively, smaller than the other experimental conditions. Expanding these subsets is a priority for future iterations.

All automated evaluation used GPT-4o as the sole judge model. Human validation of MQM rankings yielded perfect agreement ($\rho = 1.0$), and the probe designer ablation demonstrates that behavioral scores are robust to switching the probe generator from GPT-4o to Mistral Large~3. Nevertheless, cross-judge validation with alternative evaluator models would further strengthen confidence in the scoring framework. Capability question ground truth annotations are still being finalized for complete model coverage, and preliminary results should be interpreted accordingly.

\section*{Ethics Statement}

All scientific figures in \textsc{SciFig-Eval} are sourced from arXiv preprints, which are openly accessible and licensed for research use. The dataset contains no personally identifiable information. Human annotations were performed by consenting graduate researchers who were informed of the study's purpose. The benchmark is designed to reveal behavioral limitations of vision-language models under controlled conditions, not to develop adversarial attacks or to game model performance.


\bibliography{references}
\bibliographystyle{acl_natbib}

\appendix

\section{Detailed Results by Chart Type}
\label{sec:appendix-chart-type}
\begin{table*}[h]
\centering
\small
\begin{tabular}{@{}lrrr@{}}
\toprule
\textbf{Model} & \textbf{Bar} & \textbf{Line} & \textbf{Pie} \\
\midrule
GPT-5.2 & 96.2$_{\scriptstyle 95.0}^{\scriptstyle 97.3}$ & 87.8$_{\scriptstyle 85.6}^{\scriptstyle 89.9}$ & 90.0$_{\scriptstyle 87.4}^{\scriptstyle 92.4}$ \\
Gemini & 95.3$_{\scriptstyle 94.1}^{\scriptstyle 96.4}$ & 85.1$_{\scriptstyle 83.0}^{\scriptstyle 87.3}$ & 90.0$_{\scriptstyle 86.8}^{\scriptstyle 92.9}$ \\
Llama 4 & 89.5$_{\scriptstyle 87.2}^{\scriptstyle 91.6}$ & 77.7$_{\scriptstyle 75.3}^{\scriptstyle 80.1}$ & 73.0$_{\scriptstyle 67.6}^{\scriptstyle 78.1}$ \\
Qwen-235B & 88.6$_{\scriptstyle 86.5}^{\scriptstyle 90.6}$ & 77.2$_{\scriptstyle 74.6}^{\scriptstyle 79.7}$ & 73.0$_{\scriptstyle 67.0}^{\scriptstyle 78.6}$ \\
Qwen-8B & 88.9$_{\scriptstyle 86.5}^{\scriptstyle 91.1}$ & 73.7$_{\scriptstyle 71.0}^{\scriptstyle 76.3}$ & 69.9$_{\scriptstyle 63.7}^{\scriptstyle 75.9}$ \\
Qwen-30B & 83.4$_{\scriptstyle 80.7}^{\scriptstyle 86.0}$ & 71.5$_{\scriptstyle 68.6}^{\scriptstyle 74.4}$ & 62.9$_{\scriptstyle 55.8}^{\scriptstyle 69.7}$ \\
Gemma & 82.1$_{\scriptstyle 79.6}^{\scriptstyle 84.5}$ & 63.3$_{\scriptstyle 60.5}^{\scriptstyle 66.1}$ & 55.5$_{\scriptstyle 48.4}^{\scriptstyle 62.6}$ \\
Phi-4 & 72.4$_{\scriptstyle 68.8}^{\scriptstyle 75.9}$ & 58.6$_{\scriptstyle 55.0}^{\scriptstyle 62.1}$ & 49.2$_{\scriptstyle 41.4}^{\scriptstyle 57.3}$ \\
\bottomrule
\end{tabular}
\caption{Baseline MQM scores by chart type with 95\% bootstrap CI.}
\label{tab:mqm-chart-type}
\end{table*}


\section{MQM Dimension Breakdown}
\label{sec:mqm-details}
\begin{table*}[t]
\centering
\small
\begin{tabular}{@{}lrrr@{}}
\toprule
\textbf{Model} & \textbf{Accuracy} & \textbf{Compl.} & \textbf{Clarity} \\
\midrule
GPT-5.2 & 3.32$_{\scriptstyle 2.81}^{\scriptstyle 3.85}$ & 0.74$_{\scriptstyle 0.54}^{\scriptstyle 0.98}$ & 0.00$_{\scriptstyle 0.00}^{\scriptstyle 0.00}$ \\
Gemini & 4.02$_{\scriptstyle 3.48}^{\scriptstyle 4.57}$ & 0.78$_{\scriptstyle 0.58}^{\scriptstyle 0.99}$ & 0.00$_{\scriptstyle 0.00}^{\scriptstyle 0.01}$ \\
Llama 4 & 6.98$_{\scriptstyle 6.32}^{\scriptstyle 7.66}$ & 1.88$_{\scriptstyle 1.54}^{\scriptstyle 2.24}$ & 0.01$_{\scriptstyle 0.00}^{\scriptstyle 0.02}$ \\
Qwen-235B & 5.99$_{\scriptstyle 5.32}^{\scriptstyle 6.68}$ & 3.16$_{\scriptstyle 2.70}^{\scriptstyle 3.67}$ & 0.01$_{\scriptstyle 0.00}^{\scriptstyle 0.03}$ \\
Qwen-8B & 7.44$_{\scriptstyle 6.74}^{\scriptstyle 8.16}$ & 2.61$_{\scriptstyle 2.19}^{\scriptstyle 3.05}$ & 0.00$_{\scriptstyle 0.00}^{\scriptstyle 0.00}$ \\
Qwen-30B & 7.50$_{\scriptstyle 6.78}^{\scriptstyle 8.23}$ & 4.72$_{\scriptstyle 4.09}^{\scriptstyle 5.34}$ & 0.01$_{\scriptstyle 0.00}^{\scriptstyle 0.03}$ \\
Gemma & 11.52$_{\scriptstyle 10.70}^{\scriptstyle 12.33}$ & 3.18$_{\scriptstyle 2.73}^{\scriptstyle 3.64}$ & 0.03$_{\scriptstyle 0.01}^{\scriptstyle 0.06}$ \\
Phi-4 & 9.49$_{\scriptstyle 8.74}^{\scriptstyle 10.25}$ & 8.75$_{\scriptstyle 7.78}^{\scriptstyle 9.81}$ & 0.10$_{\scriptstyle 0.06}^{\scriptstyle 0.16}$ \\
\bottomrule
\end{tabular}
\caption{Mean penalty per MQM dimension (lower = fewer errors). Accuracy and Completeness weighted equally (Major=5.0, Minor=2.0). Clarity weighted lower (Major=2.5, Minor=1.0).}
\label{tab:mqm-dimensions}
\end{table*}


\section{Error Sub-Type Analysis}
\label{sec:error-subtypes}
\begin{table}[t]
\centering
\small
\begin{tabular}{@{}lr@{}}
\toprule
\textbf{Error Sub-Type} & \textbf{Count} \\
\midrule
% TODO: Populate from detailed penalty analysis \\
\bottomrule
\end{tabular}
\caption{Most common MQM error sub-types across all models and figures.}
\label{tab:error-subtypes}
\end{table}

\section{Caption Bias by Modification Type}
\label{sec:caption-bias-details}
\begin{table}[H]
\centering
\small
\begin{tabular}{@{}lrrrrr@{}}
\toprule
\textbf{Model} & \textbf{Comp. Swap} & \textbf{Rank Inv.} & \textbf{Rate Mis.} & \textbf{Trend} & \textbf{Val. Anch.} \\
\midrule
GPT-5.2 & 0.83$_{\scriptstyle 0.71}^{\scriptstyle 0.93}$ & 1.00$_{\scriptstyle 1.00}^{\scriptstyle 1.00}$ & 0.87$_{\scriptstyle 0.73}^{\scriptstyle 0.97}$ & 0.84$_{\scriptstyle 0.73}^{\scriptstyle 0.94}$ & 0.94$_{\scriptstyle 0.89}^{\scriptstyle 0.98}$ \\
Gemini & 0.93$_{\scriptstyle 0.82}^{\scriptstyle 1.00}$ & 0.91$_{\scriptstyle 0.73}^{\scriptstyle 1.00}$ & 0.87$_{\scriptstyle 0.74}^{\scriptstyle 0.97}$ & 0.80$_{\scriptstyle 0.68}^{\scriptstyle 0.90}$ & 0.95$_{\scriptstyle 0.91}^{\scriptstyle 0.98}$ \\
Llama 4 & 0.74$_{\scriptstyle 0.62}^{\scriptstyle 0.87}$ & 0.80$_{\scriptstyle 0.50}^{\scriptstyle 1.00}$ & 0.76$_{\scriptstyle 0.60}^{\scriptstyle 0.92}$ & 0.71$_{\scriptstyle 0.58}^{\scriptstyle 0.83}$ & 0.74$_{\scriptstyle 0.66}^{\scriptstyle 0.82}$ \\
Qwen-235B & 0.39$_{\scriptstyle 0.25}^{\scriptstyle 0.52}$ & 0.43$_{\scriptstyle 0.21}^{\scriptstyle 0.71}$ & 0.55$_{\scriptstyle 0.36}^{\scriptstyle 0.70}$ & 0.59$_{\scriptstyle 0.46}^{\scriptstyle 0.72}$ & 0.62$_{\scriptstyle 0.54}^{\scriptstyle 0.70}$ \\
Qwen-8B & 0.41$_{\scriptstyle 0.26}^{\scriptstyle 0.56}$ & 0.40$_{\scriptstyle 0.13}^{\scriptstyle 0.67}$ & 0.38$_{\scriptstyle 0.22}^{\scriptstyle 0.54}$ & 0.46$_{\scriptstyle 0.33}^{\scriptstyle 0.59}$ & 0.42$_{\scriptstyle 0.34}^{\scriptstyle 0.50}$ \\
Qwen-30B & 0.34$_{\scriptstyle 0.20}^{\scriptstyle 0.48}$ & 0.27$_{\scriptstyle 0.07}^{\scriptstyle 0.47}$ & 0.20$_{\scriptstyle 0.09}^{\scriptstyle 0.34}$ & 0.27$_{\scriptstyle 0.16}^{\scriptstyle 0.39}$ & 0.35$_{\scriptstyle 0.27}^{\scriptstyle 0.44}$ \\
Gemma & 0.41$_{\scriptstyle 0.25}^{\scriptstyle 0.56}$ & 0.38$_{\scriptstyle 0.15}^{\scriptstyle 0.62}$ & 0.27$_{\scriptstyle 0.12}^{\scriptstyle 0.46}$ & 0.38$_{\scriptstyle 0.26}^{\scriptstyle 0.51}$ & 0.36$_{\scriptstyle 0.28}^{\scriptstyle 0.46}$ \\
Phi-4 & 0.10$_{\scriptstyle 0.02}^{\scriptstyle 0.20}$ & 0.00$_{\scriptstyle 0.00}^{\scriptstyle 0.00}$ & 0.03$_{\scriptstyle 0.00}^{\scriptstyle 0.09}$ & 0.10$_{\scriptstyle 0.02}^{\scriptstyle 0.18}$ & 0.02$_{\scriptstyle 0.00}^{\scriptstyle 0.06}$ \\
\bottomrule
\end{tabular}
\caption{Caption bias resistance by modification type. Higher = model resisted the false claim.}
\label{tab:caption-bias-type}
\end{table}


\section{Significance Tests}
\label{sec:significance}
\begin{table*}[h]
\centering
\small
\begin{tabular}{@{}lrrr@{}}
\toprule
\textbf{Comparison} & \textbf{Diff} & \textbf{$p$} & \textbf{Sig.} \\
\midrule
\multicolumn{4}{@{}l}{\emph{Baseline MQM (vs best)}} \\
\quad gpt-5.2 vs gemini-3.1-pro & 1.4 & 0.009 & ** \\
\quad gpt-5.2 vs gemma3-27b-it & 22.5 & 0.000 & *** \\
\quad gpt-5.2 vs llama4-maverick & 10.2 & 0.000 & *** \\
\quad gpt-5.2 vs phi-4-multimodal & 29.5 & 0.000 & *** \\
\quad gpt-5.2 vs qwen3-vl-235b-a22b & 10.7 & 0.000 & *** \\
\quad gpt-5.2 vs qwen3-vl-30b-a3b & 17.1 & 0.000 & *** \\
\quad gpt-5.2 vs qwen3-vl-8b & 12.7 & 0.000 & *** \\
\midrule
\multicolumn{4}{@{}l}{\emph{Resistance (vs best)}} \\
\quad gemini-3.1-pro vs gemma3-27b-it & 0.46 & 0.000 & *** \\
\quad gemini-3.1-pro vs gpt-5.2 & 0.10 & 0.000 & *** \\
\quad gemini-3.1-pro vs llama4-maverick & 0.13 & 0.000 & *** \\
\quad gemini-3.1-pro vs phi-4-multimodal & 0.70 & 0.000 & *** \\
\quad gemini-3.1-pro vs qwen3-vl-235b-a22b & 0.16 & 0.000 & *** \\
\quad gemini-3.1-pro vs qwen3-vl-30b-a3b & 0.46 & 0.000 & *** \\
\quad gemini-3.1-pro vs qwen3-vl-8b & 0.34 & 0.000 & *** \\
\bottomrule
\end{tabular}
\caption{Paired bootstrap significance tests ($B$=10,000). Each model compared against the best. * $p<.05$, ** $p<.01$, *** $p<.001$.}
\label{tab:significance}
\end{table*}


\section{Probe Designer Ablation}
\label{sec:ablation}
\begin{table*}[t]
\centering
\small
\begin{tabular}{@{}lrrrr@{}}
\toprule
\textbf{Metric} & \textbf{GPT-4o} & \textbf{Mistral} & \textbf{$\Delta$} & \textbf{$p$} \\
\midrule
Resistance & 0.80$_{\scriptstyle 0.74}^{\scriptstyle 0.86}$ & 0.86$_{\scriptstyle 0.79}^{\scriptstyle 0.91}$ & -0.057 & 0.038 \\
Caption Bias & 0.89$_{\scriptstyle 0.83}^{\scriptstyle 0.94}$ & 0.89$_{\scriptstyle 0.82}^{\scriptstyle 0.94}$ & 0.002 & 0.489 \\
\bottomrule
\end{tabular}
\caption{Probe designer ablation. GPT-4o vs Mistral Large 3 as probe generators, GPT-5.2 as test model, GPT-4o as judge. Values show mean with 95\% bootstrap CI.}
\label{tab:ablation}
\end{table*}


\section{Stability Analysis}
\label{sec:stability-details}
\input{tables/table_a7_stability}

\section{Resistance by Chart Type}
\label{sec:resistance-chart-type}
\begin{table}[t]
\centering
\small
\begin{tabular}{@{}lrrr@{}}
\toprule
\textbf{Model} & \textbf{Bar} & \textbf{Line} & \textbf{Pie} \\
\midrule
% TODO: Compute per-chart-type resistance scores \\
\bottomrule
\end{tabular}
\caption{Hallucination resistance scores by chart type.}
\label{tab:resistance-chart-type}
\end{table}

\section{Cross-Dimensional Correlations}
\label{sec:cross-dim}
\begin{table}[t]
\centering
\footnotesize
\setlength{\tabcolsep}{2pt}
\begin{tabular}{@{}lrrrrr@{}}
\toprule
& \textbf{MQM} & \textbf{Res.} & \textbf{Cap.} & \textbf{Adm.} & \textbf{Ind.} \\
\midrule
MQM & -- & 0.95 & 0.95 & 0.83 & 0.95 \\
Resist. &  & -- & 1.00 & 0.86 & 0.93 \\
Caption &  &  & -- & 0.86 & 0.93 \\
Admit. &  &  &  & -- & 0.88 \\
Induct. &  &  &  &  & -- \\
\bottomrule
\end{tabular}
\caption{Spearman rank correlations between evaluation dimensions ($n$=8 models). Values below 0.70 suggest the dimensions capture distinct aspects of model competence.}
\label{tab:cross-dim}
\end{table}


\section{Capability Questions}
\label{sec:capability-details}
\begin{table}[t]
\centering
\small
\begin{tabular}{@{}lrrrr@{}}
\toprule
\textbf{Model} & \textbf{Count} & \textbf{Comp.} & \textbf{Compar.} & \textbf{Pattern} \\
\midrule
% TODO: Fill with capability question results \\
GPT-5.2 & -- & -- & -- & -- \\
Gemini & -- & -- & -- & -- \\
Llama 4 & -- & -- & -- & -- \\
Qwen-235B & -- & -- & -- & -- \\
Qwen-8B & -- & -- & -- & -- \\
Qwen-30B & -- & -- & -- & -- \\
Gemma & -- & -- & -- & -- \\
Phi-4 & -- & -- & -- & -- \\
\bottomrule
\end{tabular}
\caption{Capability question accuracy by category (counting, computation, comparison, pattern analysis). 250 figures, 4 questions each.}
\label{tab:capability}
\end{table}

\end{document}
